% from ega4-16-fr.tex
\oldpage[IV-4]{19}
\begin{proposition}[16.4.5]
\label{IV.16.4.5}
Supposons que $X'=X\times_S S'$, $f'$ et $u$ étant les projections canoniques.
Alors les homomorphismes canoniques $P^n(u)$ (16.4.3.3) et
$\operatorname{Gr}_1(u)$ (16.4.3.6) sont bijectifs.
\end{proposition}

On a en effet alors $X'\times_{S'}X'=(X\times_S X)\times_S S'$, et on peut donc
appliquer (16.2.3, (ii)) en remplaçant $g$ par la première projection
$p_1:X\times_S X\to X$ et $f$ par la diagonale $\Delta_f$.

On notera aussi que sous les hypothèses de (16.4.5) l'homomorphisme
$\operatorname{Gr}_\bullet(u)$ (16.4.3.5) est surjectif, mais non bijectif en
général. Toutefois (16.2.4):

\begin{corollary}[16.4.6]
\label{IV.16.4.6}
Les hypothèses étant celles de (16.4.5), supposons de plus que $w:S'\to S$ soit
plat (resp. que les $\mathscr{G}r_n(\mathscr{P}_{X/S})$ soient des
$\mathcal{O}_X$-Modules plats pour $n\leq m$); alors l'homomorphisme
\[
  \operatorname{Gr}_n(u):u^*(\mathscr{G}r_n(\mathscr{P}_{X/S}))
  \longrightarrow\mathscr{G}r_n(\mathscr{P}_{X'/S'})
\]
est bijectif pour tout $n$ (resp. pour $n\leq m$).
\end{corollary}

En effet, si $w$ est plat, il en est de même de
$v:X'\times_{S'}X'\to X\times_S X$, donc la conclusion résulte de (16.2.4).

\begin{env}[16.4.7]
\label{IV.16.4.7}
Soient $S$ un préschéma, $\mathscr{E}$ un $\mathcal{O}_S$-Module
quasi-cohérent, et posons $X=\mathbf{V}(\mathscr{E})$ (II, 1.7.8), fibré
vectoriel associé à $\mathscr{E}$, égal à
$\operatorname{Spec}(\mathscr{S}_{\mathcal{O}_S}(\mathscr{E}))$. Soit
$f:X\to S$ le morphisme structural. Pour tout ouvert $U$ de $S$, et toute
section $t\in\Gamma(U,\mathscr{E})$, $t$ s'identifie à une section de
$\mathscr{S}_{\mathcal{O}_S}(\mathscr{E})$ au-dessus de $U$; soit $t'$ son
image dans
$\Gamma(f^{-1}(U),\mathcal{O}_X)=\Gamma(U,f_*(\mathcal{O}_X))=
\Gamma(U,\mathscr{S}_{\mathcal{O}_S}(\mathscr{E}))$, et posons
\[
  \delta(t)=d_{X/S}^n(t')-t'\in
  \Gamma(f^{-1}(U),\mathscr{P}_{X/S}^n);
  \tag{16.4.7.1}
  \label{IV.16.4.7.1}
\]
il est clair que $\delta$ est un di-homomorphisme de modules (correspondant à
l'homomorphisme d'anneaux
$\Gamma(U,\mathcal{O}_S)\to\Gamma(f^{-1}(U),\mathcal{O}_X)$) de
$\Gamma(U,\mathscr{E})$ dans $\Gamma(f^{-1}(U),\mathscr{P}_{X/S}^n)$, dont
l'image appartient d'ailleurs à l'idéal d'augmentation de
$\Gamma(f^{-1}(U),\mathscr{P}_{X/S}^n)$. On en déduit (en faisant varier $U$)
un homomorphisme canonique de $\mathcal{O}_X$-Algèbres
\[
  f^*(\mathscr{S}_{\mathcal{O}_S}(\mathscr{E}))
  \longrightarrow\mathscr{P}_{X/S}^n
  \tag{16.4.7.2}
  \label{IV.16.4.7.2}
\]
et d'après la remarque précédente, si $\mathscr{K}$ est l'Idéal noyau de
l'augmentation $\mathscr{S}_{\mathcal{O}_S}(\mathscr{E})\to\mathcal{O}_S$,
l'image de $\mathscr{K}^{n+1}$ par (16.4.7.2) est nulle, si bien qu'en
factorisant par $\mathscr{K}^{n+1}$, on a finalement un homomorphisme canonique
\[
  \delta_n:f^*(\mathscr{S}_{\mathcal{O}_S}(\mathscr{E})/
  \mathscr{K}^{n+1})\longrightarrow\mathscr{P}_{X/S}^n.
  \tag{16.4.7.3}
  \label{IV.16.4.7.3}
\]
\end{env}

\begin{proposition}[16.4.8]
\label{IV.16.4.8}
Sous les conditions de (16.4.7), les homomorphismes $\delta_n$ sont bijectifs et
forment un système projectif d'isomorphismes; on en déduit un isomorphisme de
$\mathcal{O}_X$-Algèbres graduées
\[
  f^*(\mathscr{S}_{\mathcal{O}_S}^{\bullet}(\mathscr{E}))
  \longrightarrow\mathscr{G}r_\bullet(\mathscr{P}_{X/S}).
  \tag{16.4.8.1}
  \label{IV.16.4.8.1}
\]
\end{proposition}

Le fait que les homomorphismes (16.4.7.3) forment un système projectif résulte
aussitôt de leur définition. Pour prouver que ce sont des isomorphismes, il
suffira de
